\problema{Network Delay Time}{743}{src/01-grafos/743-network-delay-time.cpp}{M}

Hay \texttt{n} nodos numerados de \texttt{1} a \texttt{n}, conectados por
aristas \emph{dirigidas y con peso}: \texttt{times[i] = [u, v, w]} significa
que una señal enviada desde \texttt{u} tarda \texttt{w} unidades de tiempo en
llegar a \texttt{v}. Enviamos una señal desde el nodo \texttt{k}. Queremos el
tiempo mínimo que tarda en que la señal llegue a \emph{todos} los nodos, o
\texttt{-1} si algún nodo nunca la recibe.

\paragraph{La idea, con un ejemplo de la vida real.} Imagina que quieres
esparcir un chisme por un grupo de amigos. Le cuentas el chisme al amigo
\texttt{k}, y cada amigo, en cuanto se entera, se lo pasa a otros; pero cada
llamada tarda un rato distinto (unas personas contestan rápido, otras
tardan). Queremos saber en qué momento se entera \emph{el último} amigo,
porque hasta entonces todo el grupo ya lo sabe.

En problemas anteriores todas las conexiones ``costaban'' lo mismo (un
paso), así que bastaba avanzar por capas con BFS. Aquí cada conexión tarda
distinto, y avanzar por capas ya no sirve: el amigo que descubres primero no
siempre es al que el chisme le llega antes en tiempo real. Para eso usamos el
algoritmo de \textbf{Dijkstra}: en vez de una fila normal, usamos una
\emph{cola de prioridad} (o \emph{min-heap}), que es como una fila mágica que
siempre deja pasar primero al amigo que se va a enterar más temprano.

\begin{center}
\ttfamily
\begin{tabular}{l}
2 --1--> 1\\
2 --1--> 3 --1--> 4\\
\end{tabular}
\end{center}

\noindent Con \texttt{times = [[2,1,1],[2,3,1],[3,4,1]]}, \texttt{n = 4},
\texttt{k = 2}: desde \texttt{2} la señal llega a \texttt{1} en tiempo
\texttt{1}, a \texttt{3} en tiempo \texttt{1} y a \texttt{4} (pasando por
\texttt{3}) en tiempo \texttt{2}. El nodo más lento en enterarse es \texttt{4},
así que la respuesta es \texttt{2}.

\paragraph{Cómo lo resuelve el código, paso a paso.} Guardamos, para cada
amigo, en cuánto tiempo se entera del chisme en un arreglo \texttt{dist}. Al
principio nadie sabe nada, así que todos empiezan en ``infinito'' (un número
gigante que significa ``todavía no se entera''), salvo el amigo \texttt{k},
que arranca en $0$ porque él ya lo sabe.
\begin{enumerate}
  \item Metemos a \texttt{k} en la fila mágica (el min-heap) con tiempo $0$.
  \item Sacamos de la fila al amigo con el menor tiempo pendiente. Si el
        tiempo con el que sale ya está viejo (porque mientras tanto
        encontramos una forma más rápida de que se enterara), lo ignoramos.
  \item Si no, miramos a cada vecino al que ese amigo le puede pasar el
        chisme. Calculamos a qué hora se enteraría el vecino pasando por
        este amigo. Si es más temprano que lo que teníamos apuntado,
        actualizamos su \texttt{dist} y lo metemos a la fila. A esto se le
        llama \emph{relajar} una conexión: encontrar un camino que llega
        antes y quedarnos con él.
  \item Repetimos hasta vaciar la fila.
\end{enumerate}
¿Por qué funciona? Como ningún tiempo es negativo, la primera vez que un
amigo sale de la fila su hora ya es la definitiva: cualquier otro camino
tendría que pasar por alguien que se entera igual o más tarde, y sumarle más
tiempo solo lo empeora. (Si hubiera tiempos negativos esta garantía se
rompería, y haría falta otro método llamado Bellman-Ford.)

Al final, si algún amigo se quedó en ``infinito'' es porque el chisme nunca
le llegó: devolvemos \texttt{-1}. Si todos tienen una hora real, la respuesta
es la \emph{más grande} de todas: el grupo entero se entera justo cuando se
entera el más lento.

\lstinputlisting{src/01-grafos/743-network-delay-time.cpp}

\complejidad{$O(E \log V)$}{$O(V + E)$}

\paragraph{¿Qué significa esa complejidad?} $V$ es el número de amigos
(nodos) y $E$ el número de conexiones (aristas). $O(E \log V)$ significa que
revisamos cada conexión y, cada vez que metemos algo en la fila mágica,
mantenerla ordenada cuesta un poquito ($\log V$, que crece muy despacio). En
pocas palabras: es casi tan rápido como mirar todas las conexiones una vez,
con un pequeño costo extra por tener la fila siempre ordenada.

\paragraph{Consejos para la entrevista (Amazon).} Este es el ejemplo estrella
de Dijkstra con cola de prioridad. Suelen usarlo para ver si entiendes
\emph{por qué} el algoritmo necesita que los tiempos no sean negativos
(conviene mencionarlo aunque no te lo pregunten). Otro punto que evalúan: cómo
manejas las entradas ``viejas'' del heap. Como en la cola de prioridad de C++
no hay una forma barata de bajarle el valor a algo que ya metiste, simplemente
metemos duplicados y, al sacarlos, descartamos los que ya no coinciden con
\texttt{dist}. Casos raros: un solo amigo (\texttt{k} se entera en tiempo
\texttt{0}), un amigo al que nadie le puede llamar (nunca se entera,
\texttt{-1}), y la trampa de creer que el camino con menos saltos es siempre
el más rápido cuando en realidad hay que comparar el tiempo acumulado.
